\ifdefined\gkver\else\def\gkver{student}\fi
\documentclass[\gkver]{gaokaozhenti}
\begin{document}

\chapter{2017年4月浙江}

\begin{enumerate}
\item
%% number: 1
%% paperName: 2017年4月浙江选考第 1 题 3 分
%% typeId: 1
%% type: 单选题
%% chapter: 运动和力
%% point: 力学单位制
%% method: 无
%% score: 3
%% degree: 950
%% duplicateId: 0
%% body:
下面物理量及其对应的国际单位制单位符号，正确的是\xzanswer{D}

\fourchoices[answer=D]
{力，$\Ukg$}
{功率，$\UJ$}
{电场强度，$\mathrm{C}/\mathrm{N}$}
{电压，$\UV$}

%% answer: D
%% memo:
\memoanswer{
	力的国际单位为 $\UN$，功率的国际单位为 $\UW$，电场强度的国际单位为 $\UNC$，电压的国际单位为 $\UV$，故 D 正确。
}

\item
%% number: 2
%% paperName: 2017年4月浙江选考第 2 题 3 分
%% typeId: 1
%% type: 单选题
%% chapter: 直线运动
%% point: 路程与位移
%% method: 无
%% score: 3
%% degree: 950
%% duplicateId: 0
%% body:
下列各组物理量中均为矢量的是\xzanswer{B}

\fourchoices[answer=B]
{路程和位移}
{速度和加速度}
{力和功}
{电场强度和电势}

%% answer: B
%% memo:
\memoanswer{
	矢量是既有大小又有方向的物理量，例如：位移、速度、加速度、力、电场强度，故 B 正确。
}

\item
%% number: 3
%% paperName: 2017年4月浙江选考第 3 题 3 分
%% typeId: 1
%% type: 单选题
%% chapter: 其他
%% point: 物理学史
%% method: 无
%% score: 3
%% degree: 850
%% duplicateId: 0
%% body:
下列描述正确的是\xzanswer{A}

\fourchoices[answer=A]
{开普勒提出所有行星绕太阳运动的轨道是椭圆}
{牛顿通过实验测出了万有引力常数}
{库仑通过扭秤实验测定了电子的电荷量}
{法拉第发现了电流的磁效应}

%% answer: A
%% memo:
\memoanswer{
	开普勒提出了行星运动三定律，指出所有太阳系中的行星的运行轨道都是椭圆，故 A 正确；卡文迪许通过扭秤实验测出了万有引力常数，故 B 错误；密立根通过油滴实验测定了电子的电荷量，故 C 错误；奥斯特发现了电流的磁效应，故 D 错误。
}

\item
%% number: 4
%% paperName: 2017年4月浙江选考第 4 题 3 分
%% typeId: 1
%% type: 单选题
%% chapter: 直线运动
%% point: 自由落体运动
%% method: 无
%% score: 3
%% degree: 850
%% duplicateId: 0
%% body:
拿一个长约 $1.5\Um$ 的玻璃筒，一端封闭，另一端有开关，把金属片和小羽毛放到玻璃筒里。把玻璃筒倒立过来，观察它们下落的情况，然后把玻璃筒里的空气抽出，再把玻璃筒倒立过来，再次观察它们下落的情况，下列说法正确的是\xzanswer{C}

\onepicture[h=4.0cm]{figs/04.png}

\fourchoices[answer=C]
{玻璃筒充满空气时，金属片和小羽毛下落一样快}
{玻璃筒充满空气时，金属片和小羽毛均做自由落体运动}
{玻璃筒抽出空气后，金属片和小羽毛下落一样快}
{玻璃筒抽出空气后，金属片比小羽毛下落快}

%% answer: C
%% memo:
\memoanswer{
	抽出空气之后，羽毛和金属片下落仅受重力作用，因此加速度一样大，所以下落一样快，故 C 正确。
}

\item
%% number: 5
%% paperName: 2017年4月浙江选考第 5 题 3 分
%% typeId: 1
%% type: 单选题
%% chapter: 直线运动
%% point: 匀速直线运动
%% method: 无
%% score: 3
%% degree: 850
%% duplicateId: 0
%% body:
4 月的江南，草长莺飞，桃红柳绿，雨水连绵。伴随温柔的雨势时常出现瓢泼大雨，雷电交加的景象，在某次闪电过后约 $2\Us$ 小明听到雷声，则雷电生成处离小明的距离约为\xzanswer{A}

\fourchoices[answer=A]
{$6\times10^{2}\Um$}
{$6\times10^{4}\Um$}
{$6\times10^{6}\Um$}
{$6\times10^{8}\Um$}

%% answer: A
%% memo:
\memoanswer{
	声速大约为 $340\Ums$，所以雷电生成处离小明的距离约为 $340\Ums\times2\Us=680\Um$，故 A 正确。
}

\item
%% number: 6
%% paperName: 2017年4月浙江选考第 6 题 3 分
%% typeId: 1
%% type: 单选题
%% chapter: 直线运动
%% point: 多段匀变速直线运动
%% method: 无
%% score: 3
%% degree: 750
%% duplicateId: 0
%% body:
汽车以 $10\Ums$ 的速度在马路上匀速行驶，驾驶员发现正前方 $15\Um$ 处的斑马线上有行人，于是刹车礼让，汽车恰好停在斑马线前，假设驾驶员反应时间为 $0.5\Us$。汽车运动的 $v-t$ 图如图所示，则汽车的加速度大小为\xzanswer{C}

\onepicture[w=4.2cm]{\includesvg{figs/06}}

\fourchoices[answer=C]
{$20\Umsq$}
{$6\Umsq$}
{$5\Umsq$}
{$4\Umsq$}

%% answer: C
%% memo:
\memoanswer{
	在驾驶员反应时间内，汽车做匀速运动，位移为 $x_1=vt=10\times0.5\Um=5\Um$，所以汽车在减速阶段的位移 $x_2=x-x_1=15\Um-5\Um=10\Um$，由 $x_2=\frac{v+0}{2}t_2$ 得匀减速运动的时间 $t_2=2\Us$，所以匀减速的加速度 $a=\frac{0-10}{2}\Umsq=-5\Umsq$，故 C 正确。
}

\item
%% number: 7
%% paperName: 2017年4月浙江选考第 7 题 3 分
%% typeId: 1
%% type: 单选题
%% chapter: 力物体的平衡
%% point: 三力平衡
%% method: 无
%% score: 3
%% degree: 750
%% duplicateId: 0
%% body:
重型自卸车利用液压装置使车厢缓慢倾斜到一定角度，车厢上的石块就会自动滑下，以下说法正确的是\xzanswer{C}

\onepicture[w=4.5cm]{figs/07.jpg}

\fourchoices[answer=C]
{在石块下滑前后自卸车与石块整体的重心位置不变}
{自卸车车厢倾角越大，石块与车厢的动摩擦因数越小}
{自卸车车厢倾角变大，车厢与石块间的正压力减小}
{石块开始下滑时，受到的摩擦力大于重力沿斜面方向的分力}

%% answer: C
%% memo:
\memoanswer{
	物体重心的位置跟形状还有质量分布有关，石块下滑前后，质量分布变化，形状变化，所以重心位置改变，故 A 错误；动摩擦因数与车厢倾角无关，故 B 错误；根据平衡关系，利用正交分解法可知 $F_{N}=mg\cos\theta$，$F_{f}=mg\sin\theta$，车厢倾角变大，所以正压力随角度变大而减小，故 C 正确；石块下滑时，重力的分力大于受到的摩擦力，故 D 错误。

	\onepicture[w=4.0cm]{figs/07_an.png}
}

\item
%% number: 8
%% paperName: 2017年4月浙江选考第 8 题 3 分
%% typeId: 1
%% type: 单选题
%% chapter: 电场
%% point: 带电粒子的偏转
%% method: 无
%% score: 3
%% degree: 650
%% duplicateId: 0
%% body:
如图所示，在竖立放置间距为 $d$ 的平行板电容器中，存在电场强度为 $E$ 的匀强电场。有一质量为 $m$，电荷量为 $+q$ 的点电荷从两极板正中间处静止释放，重力加速度为 $g$。则点电荷运动到负极板的过程\xzanswer{B}

\onepicture[w=2.8cm]{figs/08.png}

\fourchoices[answer=B]
{加速度大小为 $a=\frac{qE}{m}+g$}
{所需的时间为 $t=\sqrt{\frac{md}{qE}}$}
{下降的高度为 $y=\frac{d}{2}$}
{电场力所做的功为 $W=Eqd$}

%% answer: B
%% memo:
\memoanswer{
	点电荷受到重力、电场力的作用，由力的合成和牛顿第二定律得 $a=\frac{\sqrt{(qE)^2+(mg)^2}}{m}$，故 A 错误；根据运动的独立性和等时性，设水平方向点电荷的运动时间为 $t$，则 $\frac{d}{2}=\frac{1}{2}\frac{qE}{m}t^2$，化简得 $t=\sqrt{\frac{md}{qE}}$，故 B 正确；下降高度 $y=\frac{1}{2}gt^2=\frac{mgd}{2qE}$，故 C 错误；电场力做的功 $W=\frac{qEd}{2}$，故 D 错误。

	\onepicture[w=2.8cm]{figs/08_an.jpg}
}

\item
%% number: 9
%% paperName: 2017年4月浙江选考第 9 题 3 分
%% typeId: 1
%% type: 单选题
%% chapter: 磁场
%% point: 左手定则
%% method: 无
%% score: 3
%% degree: 550
%% duplicateId: 0
%% body:
如图所示，两平行直导线 cd 和 ef 竖直放置，通以方向相反大小相等的电流，a、b 两点位于两导线所在的平面内。则\xzanswer{D}

\onepicture[h=4.0cm]{figs/09.png}

\fourchoices[answer=D]
{b 点的磁感应强度为零}
{ef 导线在 a 点产生的磁场方向垂直纸面向里}
{cd 导线受到的安培力方向向右}
{同时改变了导线的电流方向，cd 导线受到的安培力方向不变}

%% answer: D
%% memo:
\memoanswer{
	由右手定则可知导线 cd、ef 在 b 处产生磁场方向均垂直纸面向外，故 A 错误；由右手定则知，ef 在 a 处的磁场垂直纸面向外，故 B 错误；根据左手定则可判断，方向相反的两个电流受到的安培力互相排斥，故 C 错误；只要电流方向相反，就互相排斥，故 D 正确。
}

\item
%% number: 10
%% paperName: 2017年4月浙江选考第 10 题 3 分
%% typeId: 1
%% type: 单选题
%% chapter: 力物体的平衡
%% point: 三力平衡
%% method: 整体与隔离
%% score: 3
%% degree: 650
%% duplicateId: 0
%% body:
重力为 $G$ 的体操运动员在进行自由体操比赛时，有如图所示的比赛动作，当运动员竖直倒立保持静止状态时，两手臂对称支撑，夹角为 $\theta$，则\xzanswer{A}

\onepicture[w=3.0cm]{figs/10.jpg}

\fourchoices[answer=A]
{当 $\theta=60^\circ$ 时，运动员单手对地面的正压力大小为 $\frac{G}{2}$}
{当 $\theta=120^\circ$ 时，运动员单手对地面的正压力大小为 $G$}
{当 $\theta$ 不同时，运动员受到的合力不同}
{当 $\theta$ 不同时，运动员与地面之间的相互作用力不相等}

%% answer: A
%% memo:
\memoanswer{
	将运动员看作一个整体可知，不管双臂夹角多少，$2F_{N}=G$，故 A 正确 B 错误；运动员处于静止状态，其受到的合力为零，故 C 错误；不管角度如何，相互作用力总是等大反向，故 D 错误。

	\onepicture[w=3.0cm]{figs/10_an.png}
}

\item
%% number: 11
%% paperName: 2017年4月浙江选考第 11 题 3 分
%% typeId: 1
%% type: 单选题
%% chapter: 曲线运动
%% point: 天体运动
%% method: 无
%% score: 3
%% degree: 550
%% duplicateId: 0
%% body:
如图所示，设行星绕太阳的运动是匀速圆周运动，金星自身的半径是火星的 $n$ 倍，质量为火星的 $k$ 倍，不考虑行星自转的影响，则\xzanswer{B}

\onepicture[w=5.5cm]{figs/11.jpg}

\fourchoices[answer=B]
{金星表面的重力加速度是火星的 $\frac{k}{n}$}
{金星的第一宇宙速度是火星的 $\sqrt{\frac{k}{n}}$}
{金星绕太阳运动的加速度比火星小}
{金星绕太阳运动的周期比火星大}

%% answer: B
%% memo:
\memoanswer{
	由 $g=\frac{GM}{R^{2}}$ 可知 $\frac{g_{\text{金}}}{g_{\text{火}}}=\frac{k}{n^{2}}$，故 A 错误；由 $v_{1}=\sqrt{\frac{GM}{R}}$ 可知 $\frac{v_{\text{金}}}{v_{\text{火}}}=\sqrt{\frac{k}{n}}$，故 B 正确；由 $a=\frac{GM_{\text{日}}}{r^{2}}$ 可知，轨道半径越大，加速度越小，火星的轨道半径大于金星，故火星的加速度小，故 C 错误；由 $\frac{R^{3}}{T^{2}}=C$ 可知轨道半径越大，周期越长，故火星绕太阳运动的周期比金星长，故 D 错误。
}

\item
%% number: 12
%% paperName: 2017年4月浙江选考第 12 题 3 分
%% typeId: 1
%% type: 单选题
%% chapter: 运动和力
%% point: 超重失重
%% method: 无
%% score: 3
%% degree: 550
%% duplicateId: 0
%% body:
火箭发射回收是航天技术的一大进步。如图所示，火箭在返回地面前的某段运动，可看成先匀速后减速的直线运动，最后撞落在地面上。不计火星质量的变化，则\xzanswer{D}

\onepicture[w=7.0cm]{figs/12.png}

\fourchoices[answer=D]
{火箭在匀速下降过程中机械能守恒}
{火箭在减速下降过程中携带的检测仪器处于失重状态}
{火箭在减速下降过程中合力做功，等于火箭机械能的变化}
{火箭着地时，火箭对地的作用力大于自身的重力}

%% answer: D
%% memo:
\memoanswer{
	火箭匀速下降，说明阻力等于重力，阻力做负功，所以机械能不守恒，故 A 错误；在减速过程中，加速度向上，检测仪器处于超重状态，故 B 错误；火箭着地时，处于减速阶段，地面对火箭的作用力大于火箭重力，即火箭对地的作用力大于火箭自身的重力，故 D 正确；合外力做功等于动能改变量，故 C 错误。
}

\item
%% number: 13
%% paperName: 2017年4月浙江选考第 13 题 3 分
%% typeId: 1
%% type: 单选题
%% chapter: 曲线运动
%% point: 平抛运动
%% method: 无
%% score: 3
%% degree: 550
%% duplicateId: 0
%% body:
图中给出某一通关游戏的示意图，安装在轨道 AB 上可上下移动的弹射器，能水平射出速度大小可调节的弹丸，弹丸射出口在 B 点的正上方，竖直面内的半圆弧 BCD 的半径为 $R=2.0\Um$，直径 BD 水平且与轨道 AB 处在同一竖直平面内，小孔 P 和圆心 O 连线与水平方向夹角为 $37^\circ$，游戏要求弹丸垂直于 P 点圆弧切线方向射入小孔 P 就能进入下一关。为了能通关，弹射器离 B 点的高度和弹丸射出的初速度分别是（不计空气阻力）\xzanswer{A}

\onepicture[w=5.0cm]{figs/13.png}

\fourchoices[answer=A]
{$0.15\Um$，$4\sqrt{3}\Ums$}
{$1.50\Um$，$4\sqrt{3}\Ums$}
{$0.15\Um$，$2\sqrt{6}\Ums$}
{$1.50\Um$，$2\sqrt{6}\Ums$}

%% answer: A
%% memo:
\memoanswer{
	由平抛运动规律得 $x=R+R\cos37^\circ=3.6\Um$，$y=\frac{1}{2}gt^2=h+R\sin37^\circ$，$\frac{gt}{v_0}=\tan37^\circ$，由 $x=v_0t$ 得，$v_0=4\sqrt{3}\Ums$，$h=0.15\Um$，故 A 正确。
}

\item
%% number: 14
%% paperName: 2017年4月浙江选考第 14 题 2 分
%% typeId: 2
%% type: 多选题
%% chapter: 原子和原子核
%% point: 质能方程
%% method: 无
%% score: 2
%% degree: 850
%% duplicateId: 0
%% body:
下列说法正确的是\xzanswer{BC}

\fourchoices[answer=BC]
{$\beta$，$\gamma$ 射线都是电磁波}
{原子核中所有核子单独存在时，质量总和大于该原子核的总质量}
{在 LC 振荡电路中，电容器刚放电时电容器极板上电量最多，回路电流最小}
{处于 $n=4$ 激发态的氢原子，共能辐射出四种不同频率的光子}

%% answer: BC
%% memo:
\memoanswer{
	原子核的质量小于组成它的核子的质量之和，故 B 正确；LC 振荡电路中，电容器开始放电时，由于自感线圈阻碍作用，回路电流最小，之后逐渐变大，电荷量减少，故 C 正确；$\beta$ 射线为高速电子流，不是电磁波，故 A 错误；处于 $n=4$ 激发态的氢原子最多可以放出 $C_4^2=6$ 种不同频率的光子，故 D 错误。
}

\item
%% number: 15
%% paperName: 2017年4月浙江选考第 15 题 2 分
%% typeId: 2
%% type: 多选题
%% chapter: 机械波
%% point: 波的图象
%% method: 无
%% score: 2
%% degree: 550
%% duplicateId: 0
%% body:
一列向右传播的简谐横波，当波传到 $x=2.0\Um$ 处的 P 点时开始计时，该时刻波形如图所示，$t=0.9\Us$ 时，观察到质点 P 第三次到达波峰位置，下列说法正确的是\xzanswer{BCD}

\onepicture[w=6.5cm]{\includesvg{figs/15}}

\fourchoices[answer=BCD]
{波速为 $0.5\Ums$}
{经 $1.4\Us$ 质点 P 运动的路程为 $70\Ucm$}
{$t=1.6\Us$ 时，$x=4.5\Um$ 处的质点 Q 第三次到达波谷}
{与该波发生干涉的另一列简谐横波的频率一定为 $2.5\UHz$}

%% answer: BCD
%% memo:
\memoanswer{
	由同侧法可知，起振方向为竖直向上。$t=0.9\Us$ 时，P 点第三次到达波峰，即 $\left(2+\frac{1}{4}\right)T=0.9\Us$，解得 $T=0.4\Us$，所以波速 $v=\frac{\lambda}{T}=5\Ums$，故 A 错误；$1.4\Us$ 相当于 3.5 个周期，每个周期路程为 $20\Ucm$，故经 $1.4\Us$ 质点 P 运动的路程为 $3.5\times20\Ucm=70\Ucm$，故 B 正确；经过 $t=\frac{4.5-2}{5}\Us=0.5\Us$，该波恰好传播到 Q 点，再经过 $1.1\Us=2.75T$ 后，质点 Q 第三次到达波谷，故 C 正确；要发生干涉现象，另外一列波的频率一定与该波相同，即 $2.5\UHz$，故 D 正确。
}

\item
%% number: 16
%% paperName: 2017年4月浙江选考第 16 题 2 分
%% typeId: 2
%% type: 多选题
%% chapter: 光学
%% point: 光的干涉
%% method: 无
%% score: 2
%% degree: 450
%% duplicateId: 0
%% body:
图中给出了用“双缝干涉测量光的波长”实验示意图，双缝 $S_{1}$ 和 $S_{2}$ 间距为 $0.80\Umm$，双缝到屏的距离为 $0.80\Um$，波长为 $500\Unm$ 的单色平行光垂直入射到双缝 $S_{1}$ 和 $S_{2}$ 上，在屏上形成干涉条纹，中心轴线 $OO'$ 上方第一条亮纹中心位置在 $P_{1}$ 处，第三条亮纹中心位置在 $P_{2}$ 处，现有 1 号、2 号虫子分别从 $S_{1}$ 和 $S_{2}$ 出发，以相同速度沿垂直屏方向飞行，1 号虫子到达屏后，沿屏直线爬行到 $P_{1}$，2 号虫子到达屏后，沿屏直线爬行到 $P_{2}$，假定两只虫子爬行速度均为 $10^{-3}\Ums$，正确的是\xzanswer{AB}

\onepicture[w=7.0cm]{figs/16.png}

\fourchoices[answer=AB]
{1 号虫子运动路程比 2 号短}
{两只虫子运动的时间差为 $0.2\Us$}
{两只虫子运动的时间差为 $1.0\Us$}
{已知条件不够，两只虫子运动的时间差无法计算}

%% answer: AB
%% memo:
\memoanswer{
	由条纹间距公式 $\Delta x=\frac{L}{d}\lambda$ 可知，$\Delta x=0.5\Umm$，则 1 号虫子的路程 $M_1=L+\frac{d}{2}+\Delta x$，2 号虫子的路程 $M_2=L+3\Delta x-\frac{d}{2}$，则 $\Delta M=M_2-M_1=2\Delta x-d=0.2\Umm$，$\Delta t=\frac{\Delta M}{v}=0.2\Us$，故 AB 正确 CD 错误。
}

\item
%% number: 17
%% paperName: 2017年4月浙江选考第 17 题 4 分
%% typeId: 4
%% type: 实验题
%% chapter: 曲线运动
%% point: 研究平抛运动实验
%% method: 无
%% score: 4
%% degree: 650
%% duplicateId: 0
%% body:
在研究“平抛运动”实验中，

\begin{enumerate}
	\item
	图 \ref{20174月浙江17a} 是横档条卡住平抛小球，用铅笔标注小球最高点，确定平抛运动轨迹的方法，坐标原点应选小球在斜槽末端时的 \tkanswer{B}。

	（A）球心　　（B）球的上端　　（C）球的下端

	在此实验中，下列说法正确的是 \tkanswer{BD}。（多选）

	（A）斜槽轨道必须光滑　　（B）记录的点应适当多一些　　（C）用光滑曲线把所有的点连接起来　　（D）$y$ 轴的方向根据重垂线确定
	\item
	图 \ref{20174月浙江17b} 是利用图 \ref{20174月浙江17a} 装置拍摄小球做平抛运动的频闪照片，由照片可以判断实验操作错误的是 \tkanswer{C}。

	（A）释放小球时初速度不为零　　（B）释放小球的初始位置不同　　（C）斜槽末端切线不水平
	\item
	图 3 是利用稳定的细水柱显示平抛运动轨迹的装置，其中正确的是 \tkanswer{B}。
\end{enumerate}

\twopicture[numA=1, numB=2, numstyle=arabic, labelA=20174月浙江17a, labelB=20174月浙江17b, align=right, wA=5.5cm, wB=4.2cm]
{figs/17a.png}
{figs/17b.png}

\threepicture[numA=A, numB=B, numC=C, labelprefix={}, align=right, w=2.4cm]
{figs/17c1.png}
{figs/17c2.png}
{figs/17c3.png}

%% answer: (1) B，BD；(2) C；(3) B
\jdanswer{
\begin{enumerate}
	\item
	B，BD

	\item
	C

	\item
	B

\end{enumerate}
}
%% memo:
\memoanswer{
	（1）由题意知，在实验过程中，用铅笔标注了小球的最高点，为了实验准确，坐标原点应选择在小球的上端才能保证一致，故 B 正确。实验过程中，斜槽不一定光滑，只要保证从同一位置静止释放，即使轨道粗糙，摩擦力做功也是相同的，离开斜槽末端的速度就是一样的，故 A 错误；记录点适当多一些，用平滑曲线连接，偏离较远的点应舍去，故 B 正确 C 错误；$y$ 轴必须是竖直方向，即用重垂线确定，故 D 正确。

	（2）由题图可知斜槽末端切线不水平，小球做的不是平抛运动，故 C 正确。

	（3）插入瓶中的另一根吸管使得瓶内与外界气压相等，目的就是为了保证水流流速不受瓶内水面下降而减小，能保证下降到该处的一段时间内，得到稳定的细水柱，故 B 正确。
}

\item
%% number: 18
%% paperName: 2017年4月浙江选考第 18 题 4 分
%% typeId: 4
%% type: 实验题
%% chapter: 电路
%% point: 电阻定律
%% method: 无
%% score: 4
%% degree: 650
%% duplicateId: 0
%% body:
小明用电学方法测量电线的长度，首先小明测得电线铜芯的直径为 $1.00\Umm$，估计其长度不超过 $50\Um$（已知铜的电阻率为 $1.75\times10^{-8}\UOm$），现有如下实验器材：①量程为 $3\UV$、内阻约为 $3\UkO$ 的电压表；②量程为 $0.6\UA$、内阻约为 $0.1\UO$ 的电流表；③阻值为 $0\sim2\UO$ 的滑动变阻器；④内阻可忽略，输出电压为 $3\UV$ 的电源；⑤阻值为 $R_0=4.30\UO$ 的定值电阻，开关和导线若干。小明采用伏安法测量电线电阻，正确连接电路后，调节滑动变阻器，电流表的示数从 0 开始增加，当示数为 $0.5\UA$ 时，电压表示数如图 \ref{20174月浙江18a} 所示，读数为 \tkanswer{2.50}$\UV$；根据小明测量的信息，图 \ref{20174月浙江18b} 中 P 点应该 \tkanswer{接 b}（选填“接 a”、“接 b”、“接 c”或“不接”），Q 点应该 \tkanswer{接 a}（选填“接 a”、“接 b”、“接 c”或“不接”），小明测得的电线长度为 \tkanswer{31.4}$\Um$。

\twopicture[numA=1, numB=2, numstyle=arabic, labelA=20174月浙江18a, labelB=20174月浙江18b, align=right, wA=4.5cm, wB=5.2cm]
{figs/18a.png}
{figs/18b.png}

%% answer: 2.50，接 b，接 a，31.4
\jdanswer{
2.50，接 b，接 a，31.4
}
%% memo:
\memoanswer{
	电压表量程为 $0\sim3\UV$，则其读数为 $2.50\UV$；由 $R=\rho\frac{L}{S}=\rho\frac{L}{\pi r^2}$ 可估算出电线的电阻约为 $1\UO$，为小电阻，故电流表应选择外接法，即 P 点应接 b；为了更好地调节待测电阻两端的电压，滑动变阻器应采用分压式接法，即 Q 点应接 a；由 $R_x+R_0=\frac{U}{I}$，$R_x=\rho\frac{L}{S}$，$S=\frac{\pi d^2}{4}$ 联立解得 $L=31.4\Um$。
}

\item
%% number: 19
%% paperName: 2017年4月浙江选考第 19 题 8 分
%% typeId: 6
%% type: 计算题
%% chapter: 运动和力
%% point: 牛顿定律的应用
%% method: 无
%% score: 8
%% degree: 650
%% duplicateId: 0
%% body:
游船从码头沿直线行驶到湖对岸，小明对过程进行观察，记录数据如下表，

\begin{center}
\begin{tabular}{|c|c|l|}
\hline
运动过程 & 运动时间 & 运动状态 \\
\hline
匀加速运动 & $0\sim40\Us$ & 初速度 $v_0=0$；末速度 $v=4.2\Ums$ \\
\hline
匀速运动 & $40\sim640\Us$ & $v=4.2\Ums$ \\
\hline
匀减速运动 & $640\sim720\Us$ & 靠岸时的速度 $v_t=0.2\Ums$ \\
\hline
\end{tabular}
\end{center}

\begin{enumerate}
	\item
	求游船匀加速运动过程中加速度大小 $a_{1}$，及位移大小 $x_{1}$；
	\item
	若游船和游客总质量 $M=8000\Ukg$，求游船匀减速运动过程中所受合力的大小 $F$；
	\item
	求游船在整个行驶过程中的平均速度大小。
\end{enumerate}

\onepicture[align=right, w=4.0cm]{figs/19.jpg}

%% answer: (1) $a_1=0.105\Umsq$，$x_1=84\Um$；(2) $F=400\UN$；(3) $3.86\Ums$
\jdanswer{
\begin{enumerate}
	\item
	$a_1=0.105\Umsq$，$x_1=84\Um$

	\item
	$F=400\UN$

	\item
	$3.86\Ums$

\end{enumerate}
}
%% memo:
\memoanswer{
	由题意做出 $v-t$ 图像如图。

	\onepicture[w=5.0cm]{figs/19_an.jpg}

	（1）由加速度的定义式得，游船匀加速运动的加速度 $a_{1}=\frac{\Delta v}{\Delta t_{1}}=\frac{4.2-0}{40}\Umsq=0.105\Umsq$，位移 $x_{1}=\frac{1}{2}a_{1}\Delta t_{1}^{2}=\frac{1}{2}\times0.105\times40^{2}\Um=84\Um$。

	（2）由加速度的定义式得，游船减速运动过程中的加速度 $a_{2}=\frac{\Delta v}{\Delta t_{3}}=\frac{4.2-0.2}{720-640}\Umsq=0.05\Umsq$，由牛顿第二定律得 $F=Ma_{2}=8000\times0.05\UN=400\UN$。

	（3）游船行驶的总位移大小为 $x=\frac{0+v}{2}\Delta t_{1}+v\Delta t_{2}+\frac{v+v_{t}}{2}\Delta t_{3}=84\Um+4.2\times600\Um+\frac{4.2+0.2}{2}\times80\Um=2780\Um$，由平均速度的定义知 $\overline{v}=\frac{x}{t}=\frac{2780}{720}\Ums\approx3.86\Ums$。
}

\item
%% number: 20
%% paperName: 2017年4月浙江选考第 20 题 10 分
%% typeId: 6
%% type: 计算题
%% chapter: 曲线运动
%% point: 圆周运动的应用
%% method: 无
%% score: 10
%% degree: 450
%% duplicateId: 0
%% body:
图中给出一段“S”形单行盘山公路的示意图，弯道 1、弯道 2 可看作两个不同水平面上的圆弧，圆心分别为 $O_{1}$、$O_{2}$，弯道中心线半径分别为 $r_{1}=10\Um$，$r_{2}=20\Um$，弯道 2 比弯道 1 高 $h=12\Um$，有一直道与两弯道圆弧相切。质量 $m=1200\Ukg$ 的汽车通过弯道时做匀速圆周运动，路面对轮胎的最大径向静摩擦力是车重的 1.25 倍，行驶时要求汽车不打滑。（$\sin37^\circ=0.6$，$\sin53^\circ=0.8$）

\begin{enumerate}
	\item
	求汽车沿弯道 1 中心线行驶时的最大速度 $v_{1}$；
	\item
	汽车以 $v_{1}$ 进入直道，以 $P=30\UkW$ 的恒定功率直线行驶了 $t=8.0\Us$，进入弯道 2，此时速度恰为通过弯道 2 中心线的最大速度，求直道上除重力以外的阻力对汽车做的功；
	\item
	汽车从弯道 1 的 A 点进入，从同一直径上的 B 点驶离，有经验的司机会利用路面宽度，用最短时间匀速安全通过弯道，设路宽 $d=10\Um$，求此最短时间（A、B 两点都在轨道的中心线上，计算时视汽车为质点）。
\end{enumerate}

\onepicture[align=right, w=7.0cm]{figs/20.png}

%% answer: (1) $5\sqrt{5}\Ums$；(2) $-2.1\times10^{4}\UJ$；(3) $1.8\Us$
\jdanswer{
\begin{enumerate}
	\item
	$v_{1}=5\sqrt{5}\Ums$

	\item
	$W_{f}=-2.1\times10^{4}\UJ$

	\item
	$t'=1.8\Us$

\end{enumerate}
}
%% memo:
\memoanswer{
	（1）汽车沿弯道 1 行驶时，由牛顿第二定律得 $kmg=m\frac{v_{1}^{2}}{r_{1}}$，解得 $v_{1}=\sqrt{kgr_{1}}=5\sqrt{5}\Ums$。

	（2）设弯道 2 的最大速度为 $v_{2}$，由牛顿第二定律得 $kmg=m\frac{v_{2}^{2}}{r_{2}}$，解得 $v_{2}=\sqrt{kgr_{2}}=5\sqrt{10}\Ums$，汽车在直道上，由动能定理得 $Pt-mgh+W_{f}=\frac{1}{2}mv_{2}^{2}-\frac{1}{2}mv_{1}^{2}$，代入数据可得 $W_{f}=-2.1\times10^{4}\UJ$。

	（3）沿如图所示内切的路线行驶时间最短，由几何关系得 $r'^{2}=r_{1}^{2}+\left[r'-\left(r_{1}-\frac{d}{2}\right)\right]^{2}$，代入数据可得 $r'=12.5\Um$。

	\onepicture[w=4.0cm]{figs/20_an1.jpg}

	汽车沿该线路行驶的最大速度为 $v'$，由牛顿第二定律得 $kmg=m\frac{v'^2}{r'}$，解得 $v'=\sqrt{kgr'}=12.5\Ums$，由 $\sin\theta=\frac{r_{1}}{r'}=0.8$，则对应的圆心角为 $2\theta=106^\circ$，线路长度为 $s=\frac{106^\circ}{360^\circ}\times2\pi r'\approx23.1\Um$，故最短时间为 $t'=\frac{s}{v'}\approx1.8\Us$。

	\onepicture[w=5.0cm]{figs/20_an2.jpg}
}

\item
%% number: 21
%% paperName: 2017年4月浙江选考第 21 题 2 分
%% typeId: 4
%% type: 实验题
%% chapter: 电磁感应
%% point: 验证电磁感应实验
%% method: 无
%% score: 2
%% degree: 750
%% duplicateId: 0
%% body:
\begin{enumerate}
	\item
	为完成“探究变压器线圈两端的电压与匝数的关系”的实验，必须要选用的是 \tkanswer{ACE}（多选）。

	A．有闭合铁芯的原副线圈　　B．无铁芯的原副线圈

	C．交流电源　　D．直流电源

	E．多用电表（交流电压档）　　F．多用电表（交流电流档）

	用匝数 $n_{a}=60$ 匝和 $n_{b}=120$ 匝的变压器，实验测量数据如下表，

	\begin{center}
	\begin{tabular}{|c|c|c|c|c|}
	\hline
	$U_{1}/\UV$ & 1.80 & 2.80 & 3.80 & 4.90 \\
	\hline
	$U_{2}/\UV$ & 4.00 & 6.01 & 8.02 & 9.98 \\
	\hline
	\end{tabular}
	\end{center}

	根据测量数据可判断连接电源的线圈是 \tkanswer{$n_{b}$}（填 $n_{a}$ 或 $n_{b}$）。
	\item
	用如图所示的装置做“探究感应电流方向的规律”实验，磁体从靠近线圈的上方静止下落，当磁体运动到如图所示的位置时，流过线圈的感应电流方向从 \tkanswer{b 到 a}（填“a 到 b”或“b 到 a”）。在磁体穿过整个线圈的过程中，传感器显示的电流 $i$ 随时间 $t$ 的图像应该是 \tkanswer{A}。
\end{enumerate}

\onepicture[align=right, w=3.2cm]{figs/21a.png}

\fourpicture[numA=A, numB=B, numC=C, numD=D, labelprefix={}, align=right, w=3.2cm]
{figs/21b1.png}
{figs/21b2.png}
{figs/21b3.png}
{figs/21b4.png}

%% answer: (1) ACE，$n_b$；(2) b 到 a，A
\jdanswer{
\begin{enumerate}
	\item
	ACE，$n_{b}$

	\item
	b 到 a，A

\end{enumerate}
}
%% memo:
\memoanswer{
	（1）为了完成变压器的探究，需要使用交流电源和有交流电压挡的多用电表，为了让变压效果明显，则需要含有闭合铁芯的原副线圈，故 ACE 正确；由于有漏磁，所以副线圈测量电压应该小于理论变压值，即 $n_{b}$ 为输入端，连接电源。

	（2）由楞次定律及右手定则知，当磁铁自由下落至图示位置时，电流方向从 b 到 a；根据楞次定律可知，进入电流方向与离开时电流方向相反，故 C 错误；由于离开线圈时速度比进入速度快，即切割磁感线的感应电流要大，故 A 正确 BD 错误。
}

\item
%% number: 22
%% paperName: 2017年4月浙江选考第 22 题 12 分
%% typeId: 6
%% type: 计算题
%% chapter: 电磁感应
%% point: 双杆问题
%% method: 无
%% score: 12
%% degree: 450
%% duplicateId: 0
%% body:
间距为 $l$ 的两平行金属导轨由水平部分和倾斜部分平滑连接而成，如图所示，倾角为 $\theta$ 的导轨处于大小为 $B_{1}$，方向垂直导轨平面向上的匀强磁场区间 Ⅰ 中，水平导轨上的无磁场区间静止放置一质量为 $3m$ 的“联动双杆”（由两根长为 $l$ 的金属杆，cd 和 ef，用长度为 $L$ 的刚性绝缘杆连接而成），在“联动双杆”右侧存在大小为 $B_{2}$，方向垂直导轨平面向上的匀强磁场区间 Ⅱ，其长度大于 $L$，质量为 $m$，长为 $l$ 的金属杆 ab，从倾斜导轨上端释放，达到匀速后进入水平导轨（无能量损失），杆 cd 与“联动双杆”发生碰撞后杆 ab 和 cd 合在一起形成“联动三杆”，“联动三杆”继续沿水平导轨进入磁场区间 Ⅱ 并从中滑出，运动过程中，杆 ab、cd 和 ef 与导轨始终接触良好，且保持与导轨垂直。已知杆 ab、cd 和 ef 电阻均为 $R=0.02\UO$，$m=0.1\Ukg$，$l=0.5\Um$，$L=0.3\Um$，$\theta=30^\circ$，$B_{1}=0.1\UT$，$B_{2}=0.2\UT$。不计摩擦阻力和导轨电阻，忽略磁场边界效应。求：

\begin{enumerate}
	\item
	杆 ab 在倾斜导轨上匀速运动时的速度大小 $v_{0}$；
	\item
	联动三杆进入磁场区间 II 前的速度大小 $v$；
	\item
	联动三杆滑过磁场区间 II 产生的焦耳热 $Q$。
\end{enumerate}

\onepicture[align=right, w=6.5cm]{figs/22.png}

%% answer: (1) $6\Ums$；(2) $1.5\Ums$；(3) $0.25\UJ$
\jdanswer{
\begin{enumerate}
	\item
	$v_{0}=6\Ums$

	\item
	$v=1.5\Ums$

	\item
	$Q=0.25\UJ$

\end{enumerate}
}
%% memo:
\memoanswer{
	（1）杆 ab 切割产生的感应电动势为 $E=B_{1}lv_{0}$，感应电流的大小为 $I=\frac{E}{1.5R}$，杆受到的安培力大小为 $F=B_{1}Il$，杆匀速运动条件为 $\frac{B_{1}^{2}l^{2}v_{0}}{1.5R}=mg\sin\theta$，联立解得 $v_{0}=\frac{1.5mgR\sin\theta}{B_{1}^{2}l^{2}}=6\Ums$。

	（2）杆 ab 与“联动双杆”发生碰撞时，由动量守恒得 $mv_{0}=4mv$，解得 $v=\frac{v_{0}}{4}=1.5\Ums$。

	（3）进入 $B_{2}$ 磁场区时，设速度变化量为 $\Delta v$，由动量定理得 $\overline{I}\Delta t B_{2}l=-4m\Delta v$，又 $\overline{I}\Delta t=\Delta q=\frac{B_{2}lL}{1.5R}$，联立解得 $\Delta v=-\frac{B_{2}^{2}l^{2}L}{1.5R\times4m}=-0.25\Ums$，滑出 $B_{2}$ 磁场区时，同样有 $\Delta v=-\frac{B_{2}^{2}l^{2}L}{1.5R\times4m}=-0.25\Ums$，滑出 $B_{2}$ 磁场后“联动三杆”的速度为 $v'=v+2\Delta v=1.0\Ums$，由能量守恒得 $Q=\frac{1}{2}\times4m\left(v^2-v'^2\right)=0.25\UJ$。
}

\item
%% number: 23
%% paperName: 2017年4月浙江选考第 23 题 13 分
%% typeId: 6
%% type: 计算题
%% chapter: 磁场
%% point: 带电粒子在磁场中的运动
%% method: 无
%% score: 13
%% degree: 350
%% duplicateId: 0
%% body:
如图所示，在 $xOy$ 平面内，有一电子源持续不断地沿 $x$ 正方向每秒发射出 $N$ 个速率均为 $v$ 的电子，形成宽为 $2b$，在 $y$ 轴方向均匀分布且关于 $x$ 轴对称的电子流。电子流沿 $x$ 方向射入一个半径为 $R$，中心位于原点 $O$ 的圆形匀强磁场区域，磁场方向垂直 $xOy$ 平面向里，电子经过磁场偏转后均从 P 点射出，在磁场区域的正下方有一对平行于 $x$ 轴的金属平行板 K 和 A，其中 K 板与 P 点的距离为 $d$，中间开有宽度为 $2l$ 且关于 $y$ 轴对称的小孔。K 板接地，A 与 K 两板间加有正负、大小均可调的电压 $U_{AK}$，穿过 K 板小孔到达 A 板的所有电子被收集且导出，从而形成电流。已知 $b=\frac{\sqrt{3}}{2}R$，$d=l$，电子质量为 $m$，电荷量为 $e$，忽略电子间相互作用。

\begin{enumerate}
	\item
	求磁感应强度 $B$ 的大小；
	\item
	求电子从 P 点射出时与负 $y$ 轴方向的夹角 $\theta$ 的范围；
	\item
	当 $U_{AK}=0$ 时，每秒经过极板 K 上的小孔到达极板 A 的电子数；
	\item
	画出电流 $i$ 随 $U_{AK}$ 变化的关系曲线（在答题纸上的方格纸上）。
\end{enumerate}

\onepicture[align=right, w=5.0cm]{figs/23.png}

%% answer: (1) $\frac{mv}{eR}$；(2) $-60^\circ\leqslant\theta\leqslant60^\circ$；(3) $0.82N$；(4) 见解析
\jdanswer{
\begin{enumerate}
	\item
	$B=\frac{mv}{eR}$

	\item
	$-60^\circ\leqslant\theta\leqslant60^\circ$

	\item
	$n=0.82N$

	\item
	图像如解析图所示

\end{enumerate}
}
%% memo:
\memoanswer{
	（1）由题意知，电子的轨道半径 $r=R$，由洛伦兹力提供向心力得 $evB=m\frac{v^{2}}{R}$，解得 $B=\frac{mv}{eR}$。

	（2）上端电子从 P 点射出时与负 $y$ 轴有最大夹角 $\theta_{m}$，由几何关系得 $\sin\theta_{m}=\frac{b}{R}$，解得 $\theta_{m}=60^\circ$，同理下端电子从 P 点射出时与负 $y$ 轴最大夹角也为 $60^\circ$，从 P 点射出时与负 $y$ 轴范围为 $-60^\circ\leqslant\theta\leqslant60^\circ$。

	\onepicture[w=5.0cm]{figs/23_an1.jpg}

	（3）设从小孔两端进入的电子速度方向与负 $y$ 轴方向的夹角为 $\alpha$，则 $\tan\alpha=\frac{l}{d}$，解得 $\alpha=45^\circ$，能打到 A 板上的电子的区域为 $y'=2R\sin\alpha=\sqrt{2}R$，设每秒进入两极板间的电子数为 $n$，则 $\frac{n}{N}=\frac{y'}{2b}=\frac{\sqrt{6}}{3}\approx0.82$，故 $n=0.82N$。

	\onepicture[w=3.0cm]{figs/23_an2.jpg}

	（4）设遏止电压为 $U_{c}$，由动能定理得 $eU_{c}=-\frac{1}{2}mv^{2}$，解得 $U_{c}=-\frac{1}{2e}mv^{2}$，与负 $y$ 轴成 $45^\circ$ 角的电子的运动轨迹刚好与 A 板相切，其逆过程是类平抛运动，达到饱和电流所需的最小反向电压 $U'=-\frac{1}{4e}mv^{2}$，或根据（3）可得饱和电流大小为 $i_{\max}=0.82Ne$，电流 $i$ 随 $U_{AK}$ 变化的关系图像如图所示。

	\onepicture[w=7.0cm]{figs/23_an3.png}
}

\end{enumerate}

\end{document}
